%%%%%%%%%%%%%%%%
%%% exod 3 %%%
%%%%%%%%%%%%%%%%

\Chapter{3}

\Verse{1}
{
	\hE{וּמֹשֶׁה הָיָה רֹעֶה אֶת צֹאן יִתְרוֹ חֹתְנוֹ כֹּהֵן מִדְיָן וַיִּנְהַג אֶת הַצֹּאן אַחַר הַמִּדְבָּר וַיָּבֹא אֶל הַר הָאֱלֹהִים חֹרֵבָה׃}
}
{
	\eE{
		3 And Moses had been shepherding the flock of Jethro, his
		father-in-law, priest of Midian. And he drove the flock at
		the far side of the wilderness, and he came to the
		Mountain of God, to Horeb.
	}
}
{
%	\footnote{
%		to the Mountain of God, to Horeb. Horeb and Sinai are understood to be
%		two names of the same place.
%	}
}

\Verse{2}
{
	\hE{וַיֵּרָא מַלְאַךְ יְהֹוָה אֵלָיו בְּלַבַּת אֵשׁ מִתּוֹךְ הַסְּנֶה וַיַּרְא וְהִנֵּה הַסְּנֶה בֹּעֵר בָּאֵשׁ וְהַסְּנֶה אֵינֶנּוּ אֻכָּל׃}
}
{
	\eE{
		And an angel of YHWH appeared to him in a fire’s flame
		from inside a bush. And he looked, and here: the bush was
		burning in the fire, and the bush was not consumed!
	}
}
{
%	\footnote{
%		the bush was not consumed. Insofar as this phenomenon does not conform
%		to the natural course of cause and effect, YHWH is understood, once
%		again, not as a force in nature but as outside nature, manipulating it.
%		Unlike the pagan gods, YHWH causes being. The name YHWH, which is
%		revealed to Moses in this first meeting, means “He causes to be.” That
%		is, it is a third-person masculine singular causative form of the verb
%		“to be.” This fits with the biblical picture of the deity as the
%		quintessential creator God, who creates light ex nihilo—and thus can
%		cause burning without consumption—and who is not in the natural realm of
%		existence, but rather transcends it.
%	}
}

\Verse{3}
{
	\hE{וַיֹּאמֶר מֹשֶׁה אָסֻרָה נָּא וְאֶרְאֶה אֶת הַמַּרְאֶה הַגָּדֹל הַזֶּה מַדּוּעַ לֹא יִבְעַר הַסְּנֶה׃}
}
{
	\eE{
		And Moses said, “Let me turn and see this great sight. Why
		doesn’t the bush burn?”
	}
}
{
}

\Verse{4}
{
	\hE{וַיַּרְא יְהֹוָה כִּי סָר לִרְאוֹת וַיִּקְרָא אֵלָיו אֱלֹהִים מִתּוֹךְ הַסְּנֶה וַיֹּאמֶר מֹשֶׁה מֹשֶׁה וַיֹּאמֶר הִנֵּנִי׃}
}
{
	\eE{
		And YHWH saw that he turned to see. And God called to him
		from inside the bush, and He said, “Moses, Moses.” And he
		said, “I’m here.”
	}
}
{
%	\footnote{
%		God called to him. The text does not say why Moses is selected to bring
%		Israel out of Egypt. As we observed in Genesis, it rarely does state the
%		reasons that particular persons are chosen. We only know that, once
%		Moses is chosen and commissioned, virtually the entire narrative is
%		developed through him.
%	}
}

\Verse{5}
{
	\hE{וַיֹּאמֶר אַל תִּקְרַב הֲלֹם שַׁל נְעָלֶיךָ מֵעַל רַגְלֶיךָ כִּי הַמָּקוֹם אֲשֶׁר אַתָּה עוֹמֵד עָלָיו אַדְמַת קֹדֶשׁ הוּא׃}
}
{
	\eE{
		And He said, “Don’t come close here. Take off your shoes
		from your feet, because the place on which you’re
		standing: it’s holy ground.”
	}
}
{
}

\Verse{6}
{
	\hE{וַיֹּאמֶר אָנֹכִי אֱלֹהֵי אָבִיךָ אֱלֹהֵי אַבְרָהָם אֱלֹהֵי יִצְחָק וֵאלֹהֵי יַעֲקֹב וַיַּסְתֵּר מֹשֶׁה פָּנָיו כִּי יָרֵא מֵהַבִּיט אֶל הָאֱלֹהִים׃}
}
{
	\eE{
		And He said, “I’m your father’s God, Abraham’s God,
		Isaac’s God, and Jacob’s God.” And Moses hid his face,
		because he was afraid of looking at God.
	}
}
{
%	\footnote{
%		I’m your father’s God. God’s acquaintance with Moses is personal from
%		the start. God does not introduce Himself as “I am the creator of the
%		universe” or “I am the God of Israel,” but: “I’m your father’s God.”
%	}
}

\Verse{7}
{
	\hJ{וַיֹּאמֶר יְהֹוָה רָאֹה רָאִיתִי אֶת עֳנִי עַמִּי אֲשֶׁר בְּמִצְרָיִם וְאֶת צַעֲקָתָם שָׁמַעְתִּי מִפְּנֵי נֹגְשָׂיו כִּי יָדַעְתִּי אֶת מַכְאֹבָיו׃}
}
{
	\eJ{
		And YHWH said, “I’ve seen the degradation of my people who
		are in Egypt, and I’ve heard their wail on account of
		their taskmasters, because I know their pains.
	}
}
{
}

\Verse{8}
{
	\hJ{וָאֵרֵד לְהַצִּילוֹ מִיַּד מִצְרַיִם וּלְהַעֲלֹתוֹ מִן הָאָרֶץ הַהִוא אֶל אֶרֶץ טוֹבָה וּרְחָבָה אֶל אֶרֶץ זָבַת חָלָב וּדְבָשׁ אֶל מְקוֹם הַכְּנַעֲנִי וְהַחִתִּי וְהָאֱמֹרִי וְהַפְּרִזִּי וְהַחִוִּי וְהַיְבוּסִי׃}
}
{
	\eJ{
		And I’ve come down to rescue them from Egypt’s hand and to
		bring them up from that land to a good and widespread
		land, to a land flowing with milk and honey, to the place
		of the Canaanite and the Hittite and the Amorite and the
		Perizzite and the Hivite and the Jebusite.
	}
}
{
%	\footnote{
%		the place of the Canaanite … We would expect the list of the seven
%		groups (as in Deut 7:1). Why are the Girgashites left out here (and in
%		v. 17)? It is simply a scribal error. The Girgashites are in fact
%		included in the Qumran, Septuagint, and Samaritan versions of this text.
%		We should accept the fact that not every question has a deep literary or
%		moral answer to it. Sometimes the explanation is simple, mundane, and
%		human.
%	}
}

\Verse{9}
{
	\hE{וְעַתָּה הִנֵּה צַעֲקַת בְּנֵי יִשְׂרָאֵל בָּאָה אֵלָי וְגַם רָאִיתִי אֶת הַלַּחַץ אֲשֶׁר מִצְרַיִם לֹחֲצִים אֹתָם׃}
}
{
	\eE{
		“And now, here, the cry of the children of Israel has come
		to me, and also I’ve seen the oppression that Egypt is
		causing them.
	}
}
{
}

\Verse{10}
{
	\hE{וְעַתָּה לְכָה וְאֶשְׁלָחֲךָ אֶל פַּרְעֹה וְהוֹצֵא אֶת עַמִּי בְנֵי יִשְׂרָאֵל מִמִּצְרָיִם׃}
}
{
	\eE{
		And now go, and I’ll send you to Pharaoh, and bring out my
		people, the children of Israel, from Egypt.”
	}
}
{
%	\footnote{
%		And now go, and I’ll send you to Pharaoh. Until now, it has all been
%		good news to Moses: God is aware of Israel’s problems back in Egypt and
%		is going to do something about them. But Moses still has no knowledge of
%		why God is revealing this to him. Then these words come: “Go!” “I’ll
%		send you!” And then comes the command: “And [you] bring out my people!”
%		Moses might well ask, “Who am I, that I should go?” And so he does.
%	}
}

\Verse{11}
{
	\hE{וַיֹּאמֶר מֹשֶׁה אֶל הָאֱלֹהִים מִי אָנֹכִי כִּי אֵלֵךְ אֶל פַּרְעֹה וְכִי אוֹצִיא אֶת בְּנֵי יִשְׂרָאֵל מִמִּצְרָיִם׃}
}
{
	\eE{
		And Moses said to God, “Who am I that I should go to
		Pharaoh and that I should bring out the children of Israel
		from Egypt?” 3:11–12. “Who am I?” “Because I’ll be with
		you.” Some say that Moses’ first response shows his
		extreme humility. But in that case, God’s answer is
		strange: “Because I’ll be with you” is not an answer to
		the question “Who am I?” Rather, as in the case of the
		exchange with Abraham over Sodom and Gomorrah, God
		apparently answers what is in the person’s heart rather
		than what he says with his mouth. “Because I shall be with
		you” is a response to someone who is saying in his heart:
		“I’m afraid.” Moses has just been told that he must do
		this. Now God assures him that “I’ll be with you.”
	}
}
{
}

\Verse{12}
{
	\hE{וַיֹּאמֶר כִּי אֶהְיֶה עִמָּךְ וְזֶה לְּךָ הָאוֹת כִּי אָנֹכִי שְׁלַחְתִּיךָ בְּהוֹצִיאֲךָ אֶת הָעָם מִמִּצְרַיִם תַּעַבְדוּן אֶת הָאֱלֹהִים עַל הָהָר הַזֶּה׃}
}
{
	\eE{
		And He said, “Because I’ll be with you. And this is the
		sign for you that I have sent you. When you bring out the
		people from Egypt you shall serve God on this mountain.”
	}
}
{
%	\footnote{
%		this is the sign. What is the sign? Is it the bush itself? Is it the
%		fact that he hears the divine voice from the bush? Is it that he will
%		successfully “bring the people from Egypt to serve God on this
%		mountain”? Or is it the very fact that “I’ll be with you”? We get a clue
%		near the end of the Torah, in Deuteronomy (4:11). There Moses says that
%		Mount Horeb “was burning”—using the same word that is used here for the
%		burning of the bush. This suggests that the burning bush that Moses
%		experienced was already a sign of what was to come. And this in turn
%		suggests that the sign to which God refers here is the miraculous
%		burning bush itself. And it prefigures what will happen when the people
%		come back to that mountain. This point is supported later in that same
%		chapter in Deuteronomy when Moses reminds the people that at Horeb “you
%		heard His voice from inside the fire” (4:36)—which recalls how Moses
%		himself hears God call “from inside the bush” that is on fire (Exod
%		3:4).
%	}
}

\Verse{13}
{
	\hE{וַיֹּאמֶר מֹשֶׁה אֶל הָאֱלֹהִים הִנֵּה אָנֹכִי בָא אֶל בְּנֵי יִשְׂרָאֵל וְאָמַרְתִּי לָהֶם אֱלֹהֵי אֲבוֹתֵיכֶם שְׁלָחַנִי אֲלֵיכֶם וְאָמְרוּ לִי מַה שְּׁמוֹ מָה אֹמַר אֲלֵהֶם׃}
}
{
	\eE{
		And Moses said to God, “Here, I’m coming to the children
		of Israel, and I’ll say to them, ‘Your fathers’ God sent
		me to you.’ And they’ll say to me, ‘What is His name?’
		What shall I say to them?”
	}
}
{
}

\Verse{14}
{
	\hE{וַיֹּאמֶר אֱלֹהִים אֶל מֹשֶׁה אֶהְיֶה אֲשֶׁר אֶהְיֶה וַיֹּאמֶר כֹּה תֹאמַר לִבְנֵי יִשְׂרָאֵל אֶהְיֶה שְׁלָחַנִי אֲלֵיכֶם׃}
}
{
	\eE{
		And God said to Moses, “I am who I am.” And He said, “You
		shall say this to the children of Israel: ‘I Am’ has sent
		me to you.”
	}
}
{
%	\footnote{
%		I am who I am. (Or: “I shall be who I shall be.” The imperfect verb here
%		is not limited to present or future time.) This answer to Moses’ second
%		response is the first formal presentation of the divine name, revealed
%		first in the first person, ’HYH, and thereafter in the third person,
%		YHWH. It has long been uncertain why this double identification occurs
%		here. In v. 14 God tells Moses to say “’HYH sent me,” and in v. 15 God
%		tells him to say “YHWH sent me.” It appears most likely to me that the
%		former is expressed in the first person because it is the deity’s own
%		first articulation of His identity in the world. And the second is the
%		deity’s informing Moses of the name in the form in which humans will
%		know it forever after, naturally in the third person.
%	}
}

\Verse{15}
{
	\hE{וַיֹּאמֶר עוֹד אֱלֹהִים אֶל מֹשֶׁה כֹּה תֹאמַר אֶל בְּנֵי יִשְׂרָאֵל יְהֹוָה אֱלֹהֵי אֲבֹתֵיכֶם אֱלֹהֵי אַבְרָהָם אֱלֹהֵי יִצְחָק וֵאלֹהֵי יַעֲקֹב שְׁלָחַנִי אֲלֵיכֶם זֶה שְּׁמִי לְעֹלָם וְזֶה זִכְרִי לְדֹר דֹּר׃}
}
{
	\eE{
		And God said further to Moses, “You shall say this to the
		children of Israel: YHWH, your fathers’ God, Abraham’s
		God, Isaac’s God, and Jacob’s God has sent me to you. This
		is my name forever, and this is how I am to be remembered
		for generation after generation.
	}
}
{
%	\footnote{
%		YHWH. The name of God is now revealed. It is a verb. It is third person.
%		It is singular. And it is masculine. Its root meaning is “to be.” It is
%		generally understood to be a causative form. Its tense is the imperfect,
%		and it cannot be limited to a past, present, or future time. Its nearest
%		translation would be: He Causes To Be. Regarding its masculine gender,
%		we must acknowledge that all the signs indicate that biblical Israel
%		conceived of God as male. The terms for the deity are all masculine
%		words. The word that became most associated with a feminine aspect of
%		deity in later Judaism, “Shechinah,” does not occur in the Tanak.
%	}
}

\Verse{16}
{
	\hJ{לֵךְ וְאָסַפְתָּ אֶת זִקְנֵי יִשְׂרָאֵל וְאָמַרְתָּ אֲלֵהֶם יְהֹוָה אֱלֹהֵי אֲבֹתֵיכֶם נִרְאָה אֵלַי אֱלֹהֵי אַבְרָהָם יִצְחָק וְיַעֲקֹב לֵאמֹר פָּקֹד פָּקַדְתִּי אֶתְכֶם וְאֶת הֶעָשׂוּי לָכֶם בְּמִצְרָיִם׃}
}
{
	\eJ{
		“Go and gather Israel’s elders and say to them, ‘YHWH,
		your fathers’ God, has appeared to me—the God of Abraham,
		Isaac, and Jacob—saying: I have taken account of you and
		what has been done to you in Egypt,
	}
}
{
}

\Verse{17}
{
	\hJ{וָאֹמַר אַעֲלֶה אֶתְכֶם מֵעֳנִי מִצְרַיִם אֶל אֶרֶץ הַכְּנַעֲנִי וְהַחִתִּי וְהָאֱמֹרִי וְהַפְּרִזִּי וְהַחִוִּי וְהַיְבוּסִי אֶל אֶרֶץ זָבַת חָלָב וּדְבָשׁ׃}
}
{
	\eJ{
		and I say I shall bring you up from the degradation of
		Egypt to the land of the Canaanite and the Hittite and the
		Amorite and the Perizzite and the Hivite and the Jebusite,
		to a land flowing with milk and honey.
	}
}
{
}

\Verse{18}
{
	\hJ{וְשָׁמְעוּ לְקֹלֶךָ וּבָאתָ אַתָּה וְזִקְנֵי יִשְׂרָאֵל אֶל מֶלֶךְ מִצְרַיִם וַאֲמַרְתֶּם אֵלָיו יְהֹוָה אֱלֹהֵי הָעִבְרִיִּים נִקְרָה עָלֵינוּ וְעַתָּה נֵלְכָה נָּא דֶּרֶךְ שְׁלֹשֶׁת יָמִים בַּמִּדְבָּר וְנִזְבְּחָה לַיהֹוָה אֱלֹהֵינוּ׃}
}
{
	\eJ{
		And they’ll listen to your voice, and you’ll come, you and
		Israel’s elders, to the king of Egypt, and you’ll say to
		him, ‘YHWH, God of the Hebrews, has communicated with us.
		And now, let us go on a trip of three days in the
		wilderness so we may sacrifice to YHWH, our God.’
	}
}
{
}

\Verse{19}
{
	\hJ{וַאֲנִי יָדַעְתִּי כִּי לֹא יִתֵּן אֶתְכֶם מֶלֶךְ מִצְרַיִם לַהֲלֹךְ וְלֹא בְּיָד חֲזָקָה׃}
}
{
	\eJ{
		And I know that the king of Egypt won’t allow you to go,
		and not by a strong hand,
	}
}
{
}

\Verse{20}
{
	\hJ{וְשָׁלַחְתִּי אֶת יָדִי וְהִכֵּיתִי אֶת מִצְרַיִם בְּכֹל נִפְלְאֹתַי אֲשֶׁר אֶעֱשֶׂה בְּקִרְבּוֹ וְאַחֲרֵי כֵן יְשַׁלַּח אֶתְכֶם׃}
}
{
	\eJ{
		and I’ll put out my hand and strike Egypt with all my
		wonders that I’ll do among them, and after that he’ll let
		you go.
	}
}
{
}

\Verse{21}
{
	\hJ{וְנָתַתִּי אֶת חֵן הָעָם הַזֶּה בְּעֵינֵי מִצְרָיִם וְהָיָה כִּי תֵלֵכוּן לֹא תֵלְכוּ רֵיקָם׃}
}
{
	\eJ{
		And I’ll put this people’s favor in Egypt’s eyes; and it
		will be, when you go, you won’t go empty-handed.
	}
}
{
}

\Verse{22}
{
	\hJ{וְשָׁאֲלָה אִשָּׁה מִשְּׁכֶנְתָּהּ וּמִגָּרַת בֵּיתָהּ כְּלֵי כֶסֶף וּכְלֵי זָהָב וּשְׂמָלֹת וְשַׂמְתֶּם עַל בְּנֵיכֶם וְעַל בְּנֹתֵיכֶם וְנִצַּלְתֶּם אֶת מִצְרָיִם׃}
}
{
	\eJ{
		And each woman will ask for silver articles and gold
		articles and clothes from her neighbor and from anyone
		staying in her house, and you’ll put them on your sons and
		on your daughters, and you’ll despoil Egypt.”
	}
}
{
}
